\section{The result and its comparison target}
\label{sec:result}

The proposed simplification is a calculus with \textbf{three inference
families on completely raw terms}: use a retained clause, introduce a
paired formula, and discharge a witness while performing one elementary
evaluation. There is one explicit initial axiom clause, universal
reflexivity. Numerical computation is never an implicit conversion rule.
Neither full iterator normalization nor zero normalization is needed in
this primary presentation.

The construction is compared with ProveIt's \texttt{ConciseIPRA} at commit
\texttt{11814771b749}. The inspected frontier
includes the seventy-three filed packages and all six archives then present
in \texttt{docs/incoming}. In particular, the comparison includes
\texttt{witness\_discharge}, \texttt{clause\_discharge},
\texttt{deterministic\_opening}, \texttt{first\_instance},
\texttt{inferred\_witness}, \texttt{prefix\_clauses}, and
\texttt{substitutive\_discharge}. The last group matters: a proposal that
only removes the supplied equality proof or replacement term from the
earlier Name rule would already be present there.

What is new here is the fusion of the existing \emph{computation-and-equality
evaluation rule} with witness elimination. In contrast, the earlier
three-family Name presentations put both iterator equations into numerical
normal forms. The incoming raw contextual-discharge variant retains a
separate computation family. The present raw calculus has three families
while keeping ordinary, nonnormalizing substitution.

An optional companion uses zero-normal terms and the existing single
SeedEval pattern. It also has three families, and has the smaller complete
checker. The raw variant is the cleaner recommendation when the priority
is the trusted mathematical specification; the companion is preferable
when normalized checker size is decisive. Both preserve the entire
intuitionistic first-order theory, not merely its equations.

The result is a local simplification of the current presentation. It does
not prove an absolute minimum over all axiomatizations, encodings or
programming languages. Nor does it beat the repository's separate
508-byte compact codec, whose proof representation and validation boundary
are different. The whole-program comparison in Section~\ref{sec:metrics}
makes these boundaries explicit.

\section{The complete raw presentation}
\label{sec:raw}

\subsection{Syntax, clauses, and the one fixed axiom}

There is one numerical sort. Raw terms and native formulas are
\[
 t::=x\mid 0\mid S(t)\mid I(n,a;u.h),
 \qquad A::=(s=t)\mid D_{x,y}A.
\]
Only $u$ in $h$ is bound by the iterator. The two variables of $D$ are
bound in its matrix. Bound variables are renamed before capture-avoiding
substitution; the executable representation instead uses bound indices.
All finite expressions of this grammar are terms or formulas: there is
no normal-form side condition and no identification by computation.
In particular, $I(0,a;u.h)$ and $a$ are distinct terms.

A reusable clause has the form $[\vec A;y.B(y)]$, where $\vec A$ is empty
or consists of one formula, and $y$ is absent from that antecedent.
Its ordinary reading is, respectively, $\forall y\,B(y)$ or
$A\to\forall y\,B(y)$. A context is a finite list of clauses. The
abbreviation $\Gamma,A$ means addition of the unconditional constant
clause $[\,;v.A]$, with $v$ unused.

Every arithmetic proof begins with the explicitly counted clause
\[
                  \boxed{\mathcal R=[\,;v.v=v].}
\]
It supplies equality reflexivity by clause use. The native contradiction
is the formula $0=S0$. There are no other primitive axiom clauses.

\subsection{One deterministic elementary evaluation}

For arbitrary formulas $E,F$, define $\rho_E(F)$ by the following three
disjoint cases. Every match is literal syntax, including equality of the
displayed iterator bodies after the chosen representation of binders.
\begin{align}
 \rho_{s=t}\bigl(I(0,s;u.h)=r\bigr)&=(r=t),
       \tag{Z}\label{eq:rho-zero}\\
 \rho_{I(c,a;u.h)=t}\bigl(I(Sc,a;u.h)=r\bigr)
       &=\bigl(r=h[t/u]\bigr),
       \tag{S}\label{eq:rho-succ}\\
 \rho_E(F)&=F\qquad\text{in every other case}.
       \tag{Id}\label{eq:rho-id}
\end{align}
The zero case has a literal zero counter; the successor case has a
successor constructor, so they cannot overlap. In the second case the
entire predecessor iterator on the left of $E$ must match the one obtained
from the left of $F$ by deleting its counter's outer successor.
The replacement term, body, and counter are inferred from the two formulas;
no term template or computation description is separate rule data.

The substitution $h[t/u]$ is a single ordinary syntactic substitution.
The result is not normalized. An auxiliary formula $E$ with a $D$ root
always selects Id. The identity case does not derive $F$ without a premise:
the rule below explicitly requires a proof of $F$.

\subsection{All three inference families}

\paragraph{Use.}
For any retained clause $[\vec A;y.B]\in\Gamma$ and term $t$,
\[
 \frac{\{\Gamma\vdash A:A\in\vec A\}}
      {\Gamma\vdash B[t/y]}.
 \tag{Use}\label{eq:raw-use}
\]
There is no proof premise for an unconditional clause, and one for a
conditional clause. This is the only leaf mechanism; reflexivity is an
instance of Use at $\mathcal R$.

\paragraph{Introduce.}
\[
 \frac{\Gamma,B(t,\eta)\vdash B(t,Sz)}
      {\Gamma\vdash D_{x,y}B(x,y)},\qquad
 \eta=\begin{cases}
 z,&B\text{ is literally an equation and }t\equiv0,\\
 0,&\text{otherwise}.
 \end{cases}
 \tag{Intro}\label{eq:raw-intro}
\]
Here $z$ is fresh for $\Gamma$, the witness $t$, and the passive parameters
of $B$. The witness test is on raw syntax. The progress case is restricted
to literal equations; it does not introduce induction on quantified formulas.

\paragraph{Evaluate and discharge.}
For a fresh numerical parameter $b$, write
$\mathcal B_b=[B(b,0);y.B(b,y)]$. The third family is
\[
 \frac{\Gamma\vdash D_{x,y}B\qquad
       \Gamma,\mathcal B_b\vdash E\qquad
       \Gamma,\mathcal B_b\vdash F}
      {\Gamma\vdash\rho_E(F)}.
 \tag{Discharge}\label{eq:raw-discharge}
\]
The parameter $b$ is fresh for $\Gamma$ and the passive parameters of the
major formula. Both branch proofs use the same extended context. The
computed conclusion must be well scoped in the outer numerical scope;
equivalently, it has no free occurrence of $b$. The intermediate formulas
$E,F$ may contain $b$.

These are the entire primitive rules of the calculus, denoted $\mathcal E_3$.
There is no separate hypothesis, equality transport, computation,
elimination, case, induction, normalization, weakening or cut instruction.
The two nonidentity patterns inside $\rho$ remain explicitly counted in
the specification. Calling this ``three families'' does not conceal their
existence or establish a lower bound.

\section{Exact equivalence on raw open sequents}
\label{sec:raw-equivalence}

Let $\mathcal V_4$ be the same syntax, clauses, Use and Intro, with ordinary
paired elimination and the following two-case Eval family in place of
Discharge:
\[
 \frac{\Gamma\vdash s=t\quad\Gamma\vdash I(0,s;u.h)=r}
      {\Gamma\vdash r=t},\qquad
 \frac{\Gamma\vdash I(c,a;u.h)=t\quad
       \Gamma\vdash I(Sc,a;u.h)=r}
      {\Gamma\vdash r=h[t/u]}.
\]
Its ordinary paired elimination is
\[
 \frac{\Gamma\vdash D_{x,y}B\quad
       \Gamma,\mathcal B_b\vdash A}
      {\Gamma\vdash A},\qquad b\notin\operatorname{FV}(A).
\]
This is the repository's raw uniform-clause calculus: Eval is the inherited
two-case arithmetic interface, and $\mathcal R$ replaces a separate
reflexivity rule. The pinned Lean files \texttt{NativeEval.lean},
\texttt{NativeClausal.lean}, and \texttt{NativeUniformClauses.lean}
establish the corresponding reference equivalences~\cite{repo,rawformal}.
The new fusion theorem below is proved here, rather than attributed to
those existing formalizations.

\begin{lemma}[Structural operations]
\label{lem:raw-structural}
Weakening and injective renaming of free parameters are admissible in
$\mathcal E_3$ and $\mathcal V_4$.
\end{lemma}
\begin{proof}
Induct on derivations, freshening a bound eigenparameter before adding an
outer parameter or clause. Injective renaming preserves and reflects the
literal matches in $\rho$, preserves the zero/successor constructors,
commutes with capture-avoiding substitution, and preserves the Intro mode.
Weakening changes none of that data. No assertion that arbitrary numerical
substitution preserves an instruction's selected case is needed.
\end{proof}

\begin{lemma}[A compound theorem with no arithmetic premise]
\label{lem:theta}
For every clause context, $\Theta=D_{x,y}(0=0)$ is derivable in both systems,
even if that context does not contain $\mathcal R$.
\end{lemma}
\begin{proof}
Introduce at witness $0$. The temporary progress assumption and the
successor target are both $0=0$. Use that newly added unconditional
assumption. Thus the proof uses only Intro and Use.
\end{proof}

\begin{theorem}[Identity on the complete consequence relation]
\label{thm:raw-equivalence}
For every well-scoped raw clause context $\Gamma$ and formula $A$,
\[
               \Gamma\vdash_{\mathcal E_3}A
               \quad\Longleftrightarrow\quad
               \Gamma\vdash_{\mathcal V_4}A.
\]
Both translations are effective transformations of finite proof trees and
preserve the entire end sequent literally. The statement does not require
$\mathcal R$ to be present.
\end{theorem}
\begin{proof}
From left to right, translate a Discharge node's three premises. If its
$\rho$ computation selects Z or S, apply the corresponding Eval inside the
witness scope. If it selects Id, the translated branch proof of $F$ already
has the required conclusion. Apply ordinary paired elimination. Its
eigenparameter condition is exactly the outer-scope condition of Discharge.
Use and Intro translate without change.

Conversely, translate a reference elimination node by Discharge. Use as
auxiliary $E$ the same compound major $D_{x,y}B$, weakened into the fresh
branch. Since $E$ is compound, $\rho_E(A)=A$, so the original conclusion
is unchanged.

Translate either Eval node by Discharge with the fixed major proof of
$\Theta$ from Lemma~\ref{lem:theta}. Weaken its original equation premises
into that fresh scope, using them as $E,F$. The intended Z or S pattern
matches and reproduces the Eval conclusion. All data of that conclusion
are from the outer scope. Recursing over the finite tree, with fresh
parameters at every inserted scope, proves the two transformations.
\end{proof}

\paragraph{Direct ordinary soundness.}
Read $D_{x,y}B$ as
$\exists x(B(x,0)\to\forall y\,B(x,y))$ and read equality as actual
equality. Eliminate the existential at $b$. In case Z, the zero computation
equation and $s=t$ turn the second branch equation into $r=t$.
In case S, successor computation and the predecessor equation give
$r=h[t/u]$. Id simply retains the branch formula. Each conclusion can
leave the witness scope because its free occurrence of $b$ is forbidden.
This is intuitionistic reasoning and does not use a truth-table argument
or a restriction to the standard natural numbers.

\section{Where equality, computation and induction come from}
\label{sec:bootstrap}

The exact comparison theorem already transfers completeness. The following
derivations expose the arithmetic content and show that no equality law is
being silently supplied by the checker. Write
$J_s=I(0,s;u.0)$. All Eval applications below can be expanded by
Theorem~\ref{thm:raw-equivalence} into the three-family calculus.

\paragraph{Joining equations.}
From $s=t$ and reflexivity of $J_s$, Z yields $J_s=t$.
From $s=r$ and that equation, a second Z yields $t=r$.
Consequently $s=t,s=r\vdash t=r$. Reflexivity supplies symmetry, and
joining with symmetry supplies transitivity.

\paragraph{Both computation equations.}
Z on reflexivity of $a$ and of $I(0,a;u.h)$ yields
$I(0,a;u.h)=a$. S on reflexivity of the predecessor and successor
iterators yields
\[
           I(Sc,a;u.h)=h[I(c,a;u.h)/u].
\]
These are theorems. They are not reductions performed by raw substitution.

\paragraph{Transport in any term.}
Suppose $s=t$ and $h(s)=r$. Put
$J=I(0,s;u.h)$ and $K=I(S0,s;u.h)$.
Zero computation gives $J=s$; successor computation and that equation
give $K=h(s)$, using S. Joining gives $K=r$.
The zero Eval on $s=t$ and reflexivity of $J$ gives $J=t$.
The successor Eval on $J=t$ and $K=r$ now yields $r=h(t)$.
This is the inherited crossed transport rule in an arbitrary term template.
It supplies congruence and transport in an equation. Transport through a
compound formula follows by external induction on its $D$ depth, opening
the old $D$, introducing the new $D$ at the fresh nonzero-syntax witness,
and transporting its base and successor instances.

\paragraph{Equation induction.}
For a literal equation $E(y)$, choose an unused witness coordinate $x$.
The progress instance of Intro at witness $0$ turns
$\Gamma,E(z)\vdash E(Sz)$ into $D_{x,y}E(y)$. Derived ordinary elimination
opens the clause $[E(0);y.E(y)]$. A supplied base proof of $E(0)$ and Use
at any term $t$ derive $E(t)$. This is contextual induction on equations.
No induction on compound formulas has been added.

\paragraph{The case mode omitted at a literal zero witness.}
It is recovered at the raw witness $w=I(0,0;u.0)$, which is not literally
$0$ and is provably equal to $0$. Transport converts the local assumption
and successor target between the witnesses $w$ and $0$. Ordinary-assumption
cut, proved by induction on derivations after freshening parameters, inserts
the original case proof. Thus the single introduction recovers both
reference introductions. This argument uses raw syntax, and is one reason
not to silently identify computed terms in the primary calculus.

The full native logical interface and the arithmetic bridge are recorded
in Sections~\ref{sec:interpretations} and the attributed appendices. The
logical macros, predecessor construction, equality reflection and pairing
compiler are inherited; the new theorem changes their proof interface
without changing their conclusions.
